\documentclass[10pt, letterpaper]{article}

\usepackage[
    ignoreheadfoot,
    top=0.4 cm,
    bottom=0.4 cm,
    left=0.8 cm,
    right=0.8 cm,
    footskip=0.2 cm,
]{geometry}
\usepackage{titlesec}
\usepackage{tabularx}
\usepackage{array}
\usepackage[dvipsnames]{xcolor}
\definecolor{primaryColor}{RGB}{0, 0, 0}
\usepackage{enumitem}
\usepackage{amsmath}
\usepackage[
    pdftitle={Tanmay Bhuskute's CV},
    pdfauthor={Tanmay Bhuskute},
    pdfcreator={LaTeX with RenderCV},
    colorlinks=true,
    urlcolor=primaryColor
]{hyperref}
\usepackage[pscoord]{eso-pic}
\usepackage{calc}
\usepackage{bookmark}
\usepackage{lastpage}
\usepackage{changepage}
\usepackage{paracol}
\usepackage{ifthen}
\usepackage{needspace}
\usepackage{iftex}

\ifPDFTeX
    \input{glyphtounicode}
    \pdfgentounicode=1
    \usepackage[T1]{fontenc}
    \usepackage[utf8]{inputenc}
    \usepackage{lmodern}
\fi

\usepackage{XCharter}

\raggedright
\AtBeginEnvironment{adjustwidth}{\partopsep0pt}
\pagestyle{empty}
\setcounter{secnumdepth}{0}
\setlength{\parindent}{0pt}
\setlength{\topskip}{0pt}
\setlength{\columnsep}{0.15cm}
\pagenumbering{gobble}

\titleformat{\section}{\needspace{2\baselineskip}\bfseries\large}{}{0pt}{}[\vspace{1pt}\titlerule]
\titlespacing{\section}{-1pt}{0.15 cm}{0.08 cm}

\renewcommand\labelitemi{$\vcenter{\hbox{\small$\bullet$}}$}

\newenvironment{highlights}{
    \begin{itemize}[
        topsep=0.10 cm,
        parsep=0.06 cm,
        partopsep=0pt,
        itemsep=0.06 cm,
        leftmargin=12pt
    ]
}{
    \end{itemize}
}

\newenvironment{highlightsforbulletentries}{
    \begin{itemize}[
        topsep=0.10 cm,
        parsep=0.06 cm,
        partopsep=0pt,
        itemsep=0.04 cm,
        leftmargin=12pt
    ]
}{
    \end{itemize}
}

\newenvironment{onecolentry}{
    \begin{adjustwidth}{0 cm + 0.00001 cm}{0 cm + 0.00001 cm}
}{
    \end{adjustwidth}
}

\newenvironment{twocolentry}[2][]{
    \onecolentry
    \def\secondColumn{#2}
    \setcolumnwidth{\fill, 4.5 cm}
    \begin{paracol}{2}
}{
    \switchcolumn \raggedleft \secondColumn
    \end{paracol}
    \endonecolentry
}

\newenvironment{threecolentry}[3][]{
    \onecolentry
    \def\thirdColumn{#3}
    \setcolumnwidth{, \fill, 4.5 cm}
    \begin{paracol}{3}
    {\raggedright #2} \switchcolumn
}{
    \switchcolumn \raggedleft \thirdColumn
    \end{paracol}
    \endonecolentry
}

\newenvironment{header}{
    \setlength{\topsep}{0pt}\par\kern\topsep\centering\linespread{0.6}
}{
    \par\kern\topsep
}

\let\hrefWithoutArrow\href

\begin{document}
    \newcommand{\AND}{\unskip
        \cleaders\copy\ANDbox\hskip\wd\ANDbox
        \ignorespaces
    }
    \newsavebox\ANDbox
    \sbox\ANDbox{$|$}

    \begin{header}
        \fontsize{16 pt}{24 pt}\selectfont Tanmay Bhuskute

        \vspace{3 pt}

        \normalsize
        \mbox{ +1 (732) 858-3027 }%
        \kern 1.0 pt%
        \AND%
        \kern 1.0 pt%
        \mbox{ tanmay.v.bhuskute@gmail.com }%
        \kern 1.0 pt%
        \AND%
        \kern 1.0 pt%
        \mbox{ github.com/Tanmayb05 }%
        \kern 1.0 pt%
        \AND%
        \kern 1.0 pt%
        \mbox{ linkedin.com/in/tanmay-bhuskute }%
        \kern 1.0 pt%
        \AND%
        \kern 1.0 pt%
        \mbox{ San Francisco, CA }%
    \end{header}

    \vspace{2 pt}

\section{SUMMARY}
\begin{onecolentry}
Software Engineer and M.S. Computer Science candidate (Dec 2026) with production experience building backend automation, cloud platforms, Kubernetes deployment systems, and AI workflow tools using Python, FastAPI, Java, Go, AWS, GCP, Azure, Docker, Terraform, and CI/CD.
\end{onecolentry}

\section{EDUCATION}
\begin{twocolentry}{
    Aug 2024 - Dec 2026
}
    \textbf{Master of Science in Computer Science} \\
    University of Cincinnati
\end{twocolentry}
\begin{onecolentry}
Relevant Coursework: Advanced Algorithms, Deep Learning, Database Theory.
\end{onecolentry}
\vspace{0.15 cm}

\begin{twocolentry}{
    Aug 2018 - Jul 2022
}
    \textbf{Bachelor of Engineering in Computer Engineering} \\
    AISSMS College of Engineering
\end{twocolentry}
\begin{onecolentry}
Relevant Coursework: Data Mining, Machine Learning, Object-Oriented Programming.
\end{onecolentry}

\section{TECHNICAL SKILLS}
\begin{onecolentry}
\textbf{Programming:} Python, Java, Go, JavaScript, TypeScript, C/C++, SQL.
\end{onecolentry}
\vspace{0.06 cm}
\begin{onecolentry}
\textbf{Web \& Backend:} FastAPI, Node.js, Express.js, Flask, Django, REST APIs, Microservices, GraphQL.
\end{onecolentry}
\vspace{0.06 cm}
\begin{onecolentry}
\textbf{Databases:} PostgreSQL, MySQL, MongoDB, Redis, Elasticsearch, Vector Databases.
\end{onecolentry}
\vspace{0.06 cm}
\begin{onecolentry}
\textbf{Developer Tools:} Git, GitLab CI/CD, Jenkins, AWS, Azure, GCP, Docker, Kubernetes, ArgoCD, Helm, Terraform, Linux, Bash, SonarQube.
\end{onecolentry}
\vspace{0.06 cm}
\begin{onecolentry}
\textbf{AI \& Data:} Gemini, Vertex AI, LangChain, PyTorch, TensorFlow, scikit-learn, pandas, NumPy, RAG, MCP.
\end{onecolentry}

\section{WORK EXPERIENCE}
\begin{twocolentry}{Feb 2026 -- Present}
\textbf{Graduate Researcher} ---- University of Cincinnati
\end{twocolentry}
\begin{highlights}
  \item Built config driven Python LLM automation harness with Gemini, ADB, uiautomator2, and reproducible run artifacts.
  \item Demonstrated \textbf{25.7x latency reduction} across 53 runs on 20+ Android apps versus a published VLM baseline.
\end{highlights}

\vspace{0.15 cm}

\begin{twocolentry}{Jul 2022 -- Jun 2024}
\textbf{Software Development Engineer} ---- Siemens Industry Software
\end{twocolentry}
\begin{highlights}
  \item Automated \textbf{Teamcenter X} infrastructure workflows with Python Boto3, reducing deployment time by \textbf{40\%} across AWS environments.
  \item Migrated \textbf{40+ production environments} to Kubernetes using ArgoCD, Helm, GitLab CI/CD, EFS, and Vault.
  \item Designed Ansible Blue Green deployment framework, reducing production upgrade downtime by \textbf{50\%} with rollback support.
  \item Executed \textbf{44 production upgrades} in 4 months, resolving live Linux, EC2, database, and service issues.
  \item Diagnosed production infrastructure issues using CloudWatch metrics and Datadog logs across containerized services.
  \item Increased automated code coverage from \textbf{35\% to 78\%} through CI/CD quality gates and SonarQube standards.
\end{highlights}

\vspace{0.15 cm}

\begin{twocolentry}{Mar 2022 -- Jun 2022}
\textbf{Software Developer Intern} ---- Siemens Industry Software
\end{twocolentry}
\begin{highlights}
  \item Automated AWS resource provisioning with \textbf{Terraform}, reducing deployment time by \textbf{30\%}.
  \item Developed \textbf{Go + Terratest} framework, catching \textbf{95\%} of infrastructure misconfigurations before production.
\end{highlights}

\vspace{0.15 cm}

\begin{twocolentry}{Sep 2021 -- Dec 2021}
\textbf{Software Developer Intern} ---- Government of Maharashtra, Water Resource Department
\end{twocolentry}
\begin{highlights}
  \item Built full stack \textbf{React and MySQL} workflow app, digitizing records and improving operational efficiency by \textbf{40\%}.
\end{highlights}

\section{PROJECTS AND RESEARCH}
\begin{twocolentry}{Jul 2026}
\textbf{NYC Mobility Forecasting} [Python, PySpark, FastAPI, SQL, Docker]
\end{twocolentry}
\begin{highlights}
  \item Built FastAPI prediction service over a \textbf{2.3M row} dataset, improving MAPE \textbf{18.77\% to 16.52\%}.
\end{highlights}

\vspace{0.12 cm}

\begin{twocolentry}{May 2026}
\textbf{PolicyPilot Medical Benefit Policy Intelligence Portal} [React, TS, Node, Gemini, MCP]
\end{twocolentry}
\begin{highlights}
  \item Architected citation grounded RAG pipeline over Gemini, reducing hallucination risk in healthcare eligibility decisions.
  \item Designed \textbf{17 tool MCP server} for PA extraction, readiness evaluation, and cross payer comparison.
\end{highlights}

\section{EXTRACURRICULAR ACTIVITIES}
\begin{onecolentry}
\begin{highlights}
  \item \textbf{Finalist, RevolutionUC 2025 Hackathon} (300 participants, 31 schools) \,|\, \textbf{SAFe Practitioner 5.1} (Mar 2023).
\end{highlights}
\end{onecolentry}

\end{document}
